\begin{textsample}{POS Dim 1 – pe – Score 37.00 – pe/pe\_em\_32.txt}  \label{ex:f1_pos_181}
4.2 – ÁREA DO CONHECIMENTO DE LINGUAGENS E SUAS TECNOLOGIAS > “ Quem elegeu a busca não pode recusar a travessia. ” Guimarães Rosa A recente reformulação do Ensino Médio e a elaboração de \textbf{uma} Base Nacional Comum Curricular ( BNCC ) representam, no nosso entendimento, a travessia pela qual estamos passando na educação brasileira, sobretudo para os jovens a quem se destina o Ensino Médio, em busca do desenvolvimento de conhecimentos, atitudes e valores pertinentes à convivência, numa sociedade democrática e solidária.
De acordo com o referido documento, as aprendizagens essenciais estão organizadas por áreas do conhecimento, para as quais estão atribuídas competências específicas, referenciadas nas 10 competências gerais para a educação básica. \textbf{Aqui}, apresenta -se a concepção do Currículo de Pernambuco acerca da área de Linguagens e suas Tecnologias e dos componentes curriculares a ela vinculados.
A organização curricular adotada – áreas, interáreas, componentes, projetos, centros de interesse etc. – deve responder aos diferentes contextos e condições dos sistemas, das redes e das escolas de todo o país.
Para isso, é fundamental \textbf{que} a flexibilidade seja tomada como princípio obrigatório.
Independentemente da opção feita, é \textbf{preciso} destacar a necessidade de “ romper com a \textbf{centralidade} das \textbf{disciplinas} nos currículos e substituí -las por aspectos mais \textbf{globalizados} e \textbf{que} abranjam a complexidade das relações existentes entre os ramos da ciência no mundo real ” ( Parecer CNE/CEB no 5/2011 ).
Para tanto, é fundamental a adoção de tratamento \textbf{metodológico} \textbf{que} favoreça e estimule o protagonismo dos estudantes, haja vista o \textbf{que} coloca o Parecer CNE/CP no 11/2009 quando afirma \textbf{que} : > [... ] o tratamento por áreas do conhecimento não exclui necessariamente as \textbf{disciplinas}, com suas especificidades e saberes próprios \textbf{historicamente} construídos, mas, sim, \textbf{implica} o fortalecimento das relações entre elas e a sua \textbf{contextualização} para apreensão e intervenção na realidade, requerendo trabalho conjugado e cooperativo dos seus professores no planejamento e na execução dos planos de ensino. ( Parecer CNE/CP no 11/2009 ) Nesse sentido, a \textbf{contextualização}, a diversificação e a \textbf{transdisciplinaridade} são elementos fundamentais para o desenvolvimento da formação integral do estudante, consideradas as dimensões cognitivas e culturais da educação escolar, sobretudo num cenário em \textbf{que} a pluralidade das juventudes pernambucanas, urbanas e rurais, indígenas, quilombolas, ciganas e \textbf{ribeirinhas} nos desafia a dar vida ao currículo, no sentido de diminuir a distância entre o \textbf{que} se \textbf{diz} e o \textbf{que} se faz. \textbf{Ademais}, concordamos com o \textbf{que} está posto na BNCC ( 2018, p. 479 ), onde lemos \textbf{que} a adoção de tratamento \textbf{metodológico} \textbf{que} favoreça e estimule o protagonismo dos estudantes é de \textbf{suma} importância \textbf{uma} vez \textbf{que} as práticas devem evidenciar : > [... ] a \textbf{contextualização}, a diversificação e a \textbf{transdisciplinaridade} ou outras formas de interação e articulação entre diferentes campos de saberes específicos, contemplando vivências práticas e vinculando a educação escolar ao mundo do trabalho e à prática social e possibilitando o aproveitamento de estudos e o reconhecimento de saberes adquiridos nas experiências pessoais, sociais e do trabalho ( Resolução CNE/CEB no 3/2018, Art. 7, § 2o ).
Dessa forma, torna -se importante \textbf{que}, no novo cenário \textbf{que} se desenha na modernidade, voltado para a formação de \textbf{um} indivíduo crítico, competente e comprometido consigo mesmo, com o outro e com a natureza, se faça presente em \textbf{uma} educação \textbf{que} busca produzir sentidos, estimulando a compreensão da realidade \textbf{que} extrapola as barreiras disciplinares.
Os diversos saberes, \textbf{ora}, imprescindivelmente, atrelados aos contextos de sua produção, situados no tempo e no espaço, precisam integrar -se sob a perspectiva da \textbf{interdisciplinaridade}, promovendo a articulação entre os componentes de área para demonstrar \textbf{que} diferentes conhecimentos subsidiam a construção de outros, atingindo o nível da \textbf{transdisciplinaridade}, entendida \textbf{enquanto} integração contínua de conhecimentos, \textbf{um} diálogo entre campos do saber, a despeito dos componentes curriculares aos quais estejam ligados.
Quem habita o território da \textbf{interdisciplinaridade} não pode prescindir dos estudos transdisciplinares.
O cuidado construído arduamente nos dois territórios precisa ser devidamente respeitado em suas limitações, mas principalmente nas \textbf{inúmeras} possibilidades \textbf{que} se abrem para \textbf{uma} educação diferenciada onde o caráter humano se evidencia ( FAZENDA, 2012 ).
Assim sendo, compreende -se \textbf{que} as estratégias pedagógicas, \textbf{enquanto} atividades planejadas para alavancar o desempenho dos estudantes, devem priorizar temáticas \textbf{que} atendam à construção da identidade, à leitura crítica de mundo e à promoção de ações voltadas para o bem comum numa perspectiva inter e transdisciplinar. \textbf{Um} currículo \textbf{que} se proponha a integrar componentes de \textbf{uma} área, sob a perspectiva da \textbf{transdisciplinaridade}, deve permitir \textbf{uma} visão de “ todo ”, aproximada do \textbf{que} é estar no mundo e agir sobre ele, postura necessária para \textbf{que} se \textbf{consiga} conferir sentidos a \textbf{uma} escola \textbf{que} pretenda estar à altura do seu tempo. > [... ] a minha questão não é acabar com escola, é mudá -la completamente, é radicalmente fazer \textbf{que} nasça dela \textbf{um} novo ser tão \textbf{atual} quanto a tecnologia.
Eu continuo lutando no sentido de pôr a escola à altura do seu tempo.
E pôr a escola à altura do seu tempo não é soterrá -la, mas refazê -la. ( \textbf{FREIRE} \& PAPERT, 1996 ) Comprometidos, então, com a \textbf{tarefa} de refazer a escola, busca -se, nesse nível de ensino, consolidar e ampliar as aprendizagens previstas para o Ensino Fundamental, articulando esses conhecimentos às dimensões socioemocionais.
O foco da área de Linguagens e suas Tecnologias no Ensino Médio está, portanto, na ampliação da autonomia, do protagonismo e da autoria nas práticas de diferentes linguagens ( verbais, corporais e artísticas ) ; na identificação e na crítica aos diferentes usos das linguagens, \textbf{explicitando} seu poder no estabelecimento de relações ; na apreciação e na participação em diversas manifestações artísticas e culturais ; e no uso criativo das diversas mídias. ( BRASIL, 2018 ).
Para \textbf{que} isso possa ocorrer, tomamos como referência a concepção bakhtiniana de linguagem, cujo \textbf{fundamento} é o dialogismo, por acreditarmos \textbf{que}, nessa concepção, os \textbf{sujeitos}, em \textbf{um} constante processo de interação mediado pelo diálogo, se apoiam em suas relações sociais e de convivência cotidiana para formular suas falas, redigir seus textos, produzir cultura e se expressar diante do contexto social, histórico, cultural e ideológico em \textbf{que} \textbf{vivem}, elaborando, assim, discursos.
Segundo Bakhtin ( 2016, p. 113 ), “... falante e compreendedor jamais permanecem, cada \textbf{um}, em seu próprio mundo ; ao contrário, encontram -se num novo, num terceiro mundo, no mundo dos contatos ; dirigem -se \textbf{um} ao outro, entram em ativas relações \textbf{dialógicas}. ” Nesse sentido, as linguagens têm natureza social, na relação \textbf{que} estabelecem entre discurso e sociedade.
Importa ressaltar \textbf{que} os sentidos produzidos pela linguagem, representados por palavras, imagens, sons, gestos e movimentos materializam -se em cada componente curricular.
Mais \textbf{uma} vez, encontramo -nos com a perspectiva bakhtiniana, citada por Faraco ( 2003 apud Soares, 2009 ), quando entendemos \textbf{que} todos os diferentes campos da atividade humana estão ligados ao uso da linguagem.
Entendemos, então, ser esse o elo entre os componentes curriculares da área.
Conforme a BNCC, são cinco os campos de atuação social \textbf{que} servem para orientar a elaboração das habilidades específicas da área.
São eles : campo da vida pessoal, campo das práticas de estudo e pesquisa, campo jornalístico-midiático, campo de atuação na vida pública e o campo artístico.
O campo da vida pessoal busca provocar \textbf{uma} reflexão sobre as condições \textbf{que} cercam a vida \textbf{contemporânea} bem como \textbf{fomentar} nos estudantes a escolhas de estilos de vida saudáveis e sustentáveis.
Os estudos relacionados a esse campo contemplam o engajamento consciente, crítico e ético em relação às questões coletivas, por meio do mapeamento e do resgate de trajetórias, interesses, afinidades, antipatias, angústias, temores etc. \textbf{Aqui}, possibilita -se \textbf{uma} ampliação de referências e experiências culturais diversas e do conhecimento sobre si, além de abertura para experiências estéticas significativas.
O campo das práticas de estudo e pesquisa abrange o desenvolvimento do pensamento científico : a apreciação, a análise, a síntese, a reflexão, a \textbf{problematização}, a aplicação e a produção de discursos e textos expositivos, analíticos e argumentativos.
O percurso investigativo, \textbf{ora} experienciado, deve permitir além da leitura/escuta, a produção de diferentes gêneros textuais a fim de cumprir com o papel a \textbf{que} a ciência se propõe : compartilhar estudos, propondo soluções e divulgando resultados.
Em \textbf{uma} sociedade imediatista, de respostas rápidas, sem aprofundamento de perspectivas, \textbf{urge} \textbf{que} se estimule a atitude investigativa, a imersão no universo do conhecimento para \textbf{que} se \textbf{fomente} o \textbf{compromisso} com aspectos socioculturais e ambientais.
O campo jornalístico-midiático caracteriza -se pela circulação dos discursos/textos da mídia informativa e pelos discursos/textos publicitários.
Sua exploração permite construir \textbf{uma} consciência crítica e seletiva em relação à produção e circulação de informações, \textbf{posicionamentos} e induções ao consumo. \textbf{Aqui}, a área de Linguagens e suas Tecnologias se propõe a tratar temas polêmicos e \textbf{atuais} relacionados à comunicação em meio digital, a exemplo das fake news.
A educação midiática coloca a escola no \textbf{patamar} da \textbf{contemporaneidade}, trazendo os estudantes para \textbf{um} círculo de debates e vivências relacionados à cultura e ao letramento digital.
O campo de atuação na vida pública contempla os discursos/textos normativos, legais e jurídicos \textbf{que} regulam a convivência em sociedade, bem como discursos/textos propositivos e reivindicatórios ( petições, manifestos etc. ).
Sua exploração permite aos estudantes refletir e participar da vida pública, pautando -se pela ética.
E, por fim, o campo artístico se refere à circulação das manifestações artísticas em geral, contribuindo para a construção da apreciação estética significativa, a constituição das identidades, a vivência de processos criativos, o reconhecimento da diversidade, da multiculturalidade e a expressão de sentimentos e emoções.
Essa organização em torno de campos de atuação nos permite transitar por diferentes práticas de linguagens, por meio das quais os objetos de conhecimento dos diferentes componentes curriculares se relacionam de maneira \textbf{dialógica}.
Composta pelos componentes curriculares Educação Física, Arte, Língua Inglesa e Língua Portuguesa, a área de Linguagens e suas Tecnologias “ tem a responsabilidade de proporcionar oportunidades para a consolidação e a ampliação das habilidades de uso e reflexão sobre as linguagens – artísticas, corporais e verbais ( oral, ou visual-motora, como libras e escrita ) ”.
Neste sentido, as práticas socioculturais expressas por várias linguagens, elaboradas em espaços de interação juvenil, constituem \textbf{um} acervo de experiências \textbf{que} serão acionadas diante de situações-problema \textbf{que}, combinadas, podem produzir diferentes respostas ( BRASIL, 2017, p. 482 ).

% matched lemmas: ademais, aqui, atual, centralidade, compromisso, conseguir, contemporaneidade, contemporâneo, contextualização, dialógico, disciplina, dizer, enquanto, explicitar, fomentar, freire, fundamento, globalizado, historicamente, implicar, interdisciplinaridade, inúmero, metodológico, ora, patamar, posicionamento, preciso, problematização, que, ribeirinho, sujeito, sumo, tarefa, transdisciplinaridade, um, urgir, viver
\end{textsample}
